\documentclass[11pt, letterpaper]{article}
\usepackage[margin=0.8in, top=0.5in, bottom=0.8in]{geometry}
\usepackage[utf8]{inputenc}
\usepackage[T1]{fontenc}
\usepackage{hyperref}
\usepackage{enumitem}
\usepackage{charter} % Professional font (already confirmed present)

% --- Formatting ---
\hypersetup{
    colorlinks=true,
    linkcolor=blue,
    urlcolor=blue,
    pdftitle={CV - Yichao Jin},
}

% Manual section styling for robustness
\newcommand{\ressection}[1]{
    \vspace{10pt}
    \noindent\textbf{\large \textsc{#1}}
    \vspace{-2pt}
    \hrule
    \vspace{5pt}
}

\setlist[itemize]{leftmargin=1.5em, noitemsep, topsep=2pt, parsep=2pt}
\setlist[enumerate]{leftmargin=1.5em}

\begin{document}

% --- HEADER ---
\begin{center}
    {\huge \textbf{Yichao Jin}} \\
    \vspace{0.6em}
    School of Economic, Political and Policy Sciences $\cdot$ University of Texas at Dallas \\
    \vspace{0.4em}
    \href{mailto:Yichao.Jin@UTDallas.edu}{Yichao.Jin@UTDallas.edu} \quad | \quad \href{https://yichao2022.github.io}{yichao2022.github.io} \quad | \quad +1 (682) 266-4185
\end{center}

\vspace{0.2em}

% --- RESEARCH INTERESTS (Opening) ---
\ressection{Research Profile}
\noindent \textbf{Fields:} Health Economics, Behavioral Economics, Public Policy Analysis, Technology Governance. \\
\textbf{Methods:} Discrete Choice Experiments (DCE), Structural Modeling, Applied Econometrics (Causal Inference), Policy Simulation (SEIR), Machine Learning.

% --- EDUCATION ---
\ressection{Education}
\noindent \textbf{University of Texas at Dallas} \hfill Expected May 2026 \\
\textbf{Ph.D. in Public Policy and Political Economy} \\
- \textit{Dissertation:} ``Dynamics of Intertemporal Choice in Vaccination Timing: A DCE Approach'' \\
- \textit{Committee:} Dohyeong Kim (Chair), Richard Scotch, Soyoung Kwon. \\
- \textit{Distinction:} Dean of Graduate Education Dissertation Research Award (2025).

\vspace{0.4em}
\noindent \textbf{University of Texas at Dallas} \hfill Expected May 2026 \\
\textbf{M.S. in Social Data Analytics and Research}

\vspace{0.4em}
\noindent \textbf{University of Queensland} \hfill 2019 \\
\textbf{Master of Development Economics}

\vspace{0.4em}
\noindent \textbf{University of California, Riverside} \hfill 2017 \\
\textbf{B.A. in Economics}

% --- WORKING PAPERS (Core Research) ---
\ressection{Working Papers}
\begin{enumerate}
    \item \textbf{``How Waiting Time Shapes Preventive Health Behavior: Evidence from Vaccination Decisions''} (with Dohyeong Kim) \\
    \textit{Job Market Paper} \quad | \quad \textit{Under Review at Journal of Behavioral Medicine} \\
    \href{https://yichao2022.github.io/jmp.pdf}{\textbf{Link to Draft}} \quad | \quad \href{https://github.com/yichao2022/jmp_repo_draft}{\textbf{Research Repo}}
    \begin{itemize}
        \item \textbf{Research Question}: How do administrative delays (waiting time) and institutional trust interact to create behavioral barriers to preventive health decisions?
        \item \textbf{Approach}: Recovered formal time-preference parameters (hyperbolic and exponential discount rates) for delayed protection using a structural Discrete Choice Experiment (DCE, N=1,027).
        \item \textbf{Main Findings}: Identified that individuals with lower institutional trust are significantly more sensitive to delay friction. Heterogeneity analysis reveals that wait-time costs are disproportionately borne by rural and less-educated populations.
        \item \textbf{Policy Relevance}: Demonstrates that reducing administrative delay is as critical for health equity as financial incentives, providing a behavioral roadmap for trust-deficient populations.
    \end{itemize}

    \item \textbf{``How CEO Succession Influences Firm AI Innovation: The Role of Academic Experience''} (with Zhimin Tian and Yu Xiang) \\
    \textit{Under Review at Journal of Business Research}
    \begin{itemize}
        \item \textbf{Research Question}: Does the academic background of incoming CEOs drive AI-related innovation in high-tech firms?
        \item \textbf{Approach}: Analyzed 1,776 CEO succession events to identify shifts in innovation trajectory linked to academic credentials and experience.
    \end{itemize}

    \item \textbf{``The Price of Waiting: Evidence on Cash--Time Trade-Offs in Vaccination Time Discount''} \\
    \textit{In Preparation} --- Targeting \textit{Social Science \& Medicine}.
\end{enumerate}

% --- RESEARCH EXPERIENCE ---
\ressection{Professional Experience}
\noindent \textbf{Doctoral Researcher}, Spatial Health AI Research Partnership (SHARP) \hfill 2022 -- Present \\
\textit{UT Dallas, School of Economic, Political and Policy Sciences}
\begin{itemize}
    \item Lead researcher for interdisciplinary projects at the intersection of AI, geospatial analytics, and health policy.
    \item Developed large-scale computational simulations (R, Python) to evaluate pandemic preparedness.
\end{itemize}

\vspace{0.4em}
\noindent \textbf{Graduate Researcher}, Behavioral Health Economics Laboratory \hfill 2021 -- Present \\
\textit{UT Dallas}
\begin{itemize}
    \item Designed and executed large-scale behavioral experiments (N=1,027) analyzing risk-weighted health decisions and institutional trust.
\end{itemize}

% --- TEACHING ---
\ressection{Teaching}
\noindent \textbf{Teaching Assistant}, University of Texas at Dallas \hfill 2021 -- Present
\begin{itemize}
    \item \textbf{Quantitative Methods for Policy Analysis (Graduate)}: Led Econometrics labs; specialized in R programming, regression analysis, and causal inference for policy data.
    \item \textbf{Health Economics \& Public Policy (Undergraduate)}: Focused on health market failures and the economics of preventive care.
    \item \textbf{Public Policy Analysis (Undergraduate)}: Mentored students in drafting evidence-based policy memos and program evaluation frameworks.
\end{itemize}

% --- PRESENTATIONS & AWARDS ---
\ressection{Conferences \& Honors}
\begin{itemize}
    \item \textbf{APPAM Fall Research Conference} (2025): Estimating Time Discount Rate for Vaccination.
    \item \textbf{UT Dallas Graduate Research Symposium} (2025): Modeling Vaccination Decisions with Hyperbolic Discounting.
    \item \textbf{Dean of Graduate Education Dissertation Research Award} (2025), UT Dallas.
    \item \textbf{Betty \& Gifford Johnson Graduate Travel Award} (2025), UT Dallas.
    \item \textbf{Omicron Delta Epsilon}, International Honor Society for Economics.
\end{itemize}

% --- SKILLS ---
\ressection{Skills \& Certifications}
\noindent \textbf{Econometrics}: Mixed Logit (MXL), Latent Class Models (LCM), Structural Modeling, Causal Inference. \\
\textbf{Computing}: R, Stata (Advanced), Python, MATLAB, \LaTeX, SQL, GIS. \\
\textbf{Certificates}: Stanford Graduate Certificate in AI; CITI Human Subjects Protection.

% --- REFERENCES ---
\ressection{References}
\noindent \textbf{Dohyeong Kim (Chair)} \hfill \textbf{Richard Scotch} \hfill \textbf{Soyoung Kwon} \\
University of Texas at Dallas \hfill University of Texas at Dallas \hfill University of Texas at Dallas \\
\href{mailto:dohyeong.kim@utdallas.edu}{dohyeong.kim@utdallas.edu} \hfill \href{mailto:richard.scotch@utdallas.edu}{richard.scotch@utdallas.edu} \hfill \href{mailto:soyoung.kwon@utdallas.edu}{soyoung.kwon@utdallas.edu}

\vfill
\begin{center}
    \footnotesize \textit{Last Updated: March 17, 2026}
\end{center}

\end{document}
